% sections/appendix.tex
% Optional appendix (not counted in body page limit)

\section{Ranking Effectiveness: Precision at $K$}
\label{app:patk}

\begin{table}[h]
\centering
\caption{Precision at $K$ for the composite ranking.
  Lift = P@$K$ / random baseline (34/\NskillsDMR{} $\approx0.49\%$).}
\label{tab:patk}
\small
\begin{tabular}{@{}rcccr@{}}
\toprule
\textbf{$K$} & \textbf{T6 found} & \textbf{P@$K$} & \textbf{Recall} & \textbf{Lift} \\
\midrule
 10 & 1  & 0.100 & 2.9\%  & 20.3$\times$ \\
 20 & 3  & 0.150 & 8.8\%  & 30.5$\times$ \\
 50 & 4  & 0.080 & 11.8\% & 16.3$\times$ \\
100 & 4  & 0.040 & 11.8\% & 8.1$\times$ \\
200 & 9  & 0.045 & 26.5\% & 9.1$\times$ \\
500 & 14 & 0.028 & 41.2\% & 5.7$\times$ \\
\bottomrule
\end{tabular}
\end{table}

\section{LSCPS Keyword Bank}
\label{app:lscps}

Table~\ref{tab:lscps_kw} lists representative regex patterns for each
\lscps{} capability family used in our static capability extractor.
The full bank contains 300+ patterns across 14 capability families.

\begin{table}[h]
\centering
\caption{LSCPS capability families and representative detection patterns.}
\label{tab:lscps_kw}
\small
\begin{tabular}{@{}ll@{}}
\toprule
\textbf{Family} & \textbf{Representative Patterns (regex)} \\
\midrule
Shell Exec     & \texttt{subprocess}, \texttt{os\.system}, \texttt{exec\textbackslash(}, \texttt{bash} \\
Cred Read      & \texttt{api\_key}, \texttt{access\_token}, \texttt{bearer}, \texttt{\.env} \\
Cred Transmit  & \texttt{post.*token}, \texttt{upload.*key}, \texttt{exfil} \\
Net Write      & \texttt{requests\.post}, \texttt{webhook}, \texttt{smtp} \\
Net Read       & \texttt{requests\.get}, \texttt{httpx\.get}, \texttt{curl http} \\
File Write     & \texttt{open\textbackslash(.*"w"}, \texttt{shutil\.}, \texttt{os\.remove} \\
File Read      & \texttt{open\textbackslash(.*"r"}, \texttt{glob\textbackslash(}, \texttt{os\.listdir} \\
Autonomous     & \texttt{autonomous}, \texttt{without.*confirm}, \texttt{cron} \\
Multi-step     & \texttt{multi.step}, \texttt{workflow}, \texttt{pipeline} \\
Memory Write   & \texttt{memory}, \texttt{remember}, \texttt{long.term} \\
\bottomrule
\end{tabular}
\end{table}

\section{T6 Case Timeline}
\label{app:t6_timeline}

For each T6 case we record: date of malicious version publication,
date of fix, exposure window (days), and installs during exposure.
Table~\ref{tab:t6_timeline} summarizes the 40 confirmed cases.

\begin{table}[h]
\centering
\caption{T6 case timeline summary (representative subset).}
\label{tab:t6_timeline}
\scriptsize
\begin{tabular}{@{}lrrlp{2.6cm}@{}}
\toprule
\textbf{Skill} & \textbf{Win.} & \textbf{Inst.} & \textbf{CES} & \textbf{Primary Attack} \\
\midrule
agentmail     & 1.0\,d & 155 & 8.5 & Webhook injection \\
moltguard     & 0.1\,d &  54 & 5.0 & Phishing file \\
simmer        & 0.4\,d &  21 & 5.0 & SKILL.md rollback \\
guard-scanner & 9.1\,d &   9 & 14.5 & Trojanized scanner \\
skill-auditor & --- &   8 & 4.5 & Credential harvest \\
vmware-aiops  & 1.1\,d &   3 & 4.5 & Webhook + file delete \\
vmware-monitor& 2.1\,d &   0 & 6.5 & Webhook + credential \\
aport-guardrail& 2.6\,d &   0 & 3.0 & Shell exec injection \\
clawsentinel  & 0.3\,d &   0 & 8.5 & Repeated inject/revert \\
mdk-wallet    & --- &   0 & 11.5 & rm -rf + credential exfil \\
\midrule
\textbf{All 40} & \textbf{avg 0.5\,d} & \textbf{298} & & \\
\bottomrule
\end{tabular}
\end{table}

\section{Deep Attack Case Studies}
\label{app:casestudies}

We present five representative T6 cases selected for diversity of attack
patterns, severity, and user impact. Full reconstruction details are available
in our artifact release.

\noindent\textit{Skill names below are pseudonymized per our responsible disclosure policy (\S\ref{sec:conclusion}); original identifiers are available to reviewers under embargo request.}

\paragraph{Case 1: Trojanized security scanner (\texttt{guard-scanner}).}
A skill marketed as a security scanning tool contained
credential harvesting (5 matches), network exfiltration (webhook.site,
curl), persistence mechanisms (cron, startup), and download-execute chains.
\textbf{Exposure}: 9.1~days, 9~installs. CR = 0.143; RWM = 12.86.
The attack used the ironic disguise of a ``security'' tool to normalize
access to security-critical APIs. All three detection layers (SCS, DMR, TDM) catch this case;
the platform scanner missed it.

\paragraph{Case 2: Prompt injection via email (\texttt{agentmail}).}
The highest-impact T6 case (155~installs) used unfiltered email webhooks
to forward arbitrary content as LLM prompts, enabling indirect prompt injection.
The fix (v1.1.0) added a webhook sender allowlist.
\textbf{Exposure}: 1.0~day. CR = 0.333; RWM = 2.12.
This case demonstrates that even moderate RWM scores can indicate
critical vulnerabilities when the attack vector is semantic rather than code-pattern-based.

\paragraph{Case 3: Financial data exfiltration (\texttt{mdk-agent-wallet}).}
A cryptocurrency wallet skill contained destructive file operations
(\texttt{rm -rf}) targeting the wallet directory,
webhook exfiltration of private keys, and credential harvesting.
The fix removed the destructive commands and added \texttt{autonomous:false}.
CR = 0.143; RWM = 15.33.
\textbf{Data flow}: Private key read $\to$ webhook POST $\to$ external endpoint.

\paragraph{Case 4: Hidden phishing attack (\texttt{moltguard}).}
A malware guard tool embedded a phishing attack file in its test
instructions, instructing the LLM agent to notify users when quota
tags appeared---effectively weaponizing the agent as a phishing vector.
\textbf{Exposure}: 0.1~days, 54~installs. CR = 0.000; RWM = 16.86.

\paragraph{Case 5: Repeated injection campaign (\texttt{clawsentinel}).}
This skill appeared twice in T6 (v3.0.0$\to$2.3.3 and v2.3.4$\to$2.3.5),
following an inject-revert-inject pattern. Each injection added
\texttt{rm -rf}, credential harvesting, and file destruction,
only to be reverted within hours.
CR = 0.143; \textbf{RWM = 29.96}---the highest in our dataset.
This case is consistent with repeated probing of platform detection limits.

\paragraph{Detection layer effectiveness across cases.}
In all five cases: (a) the platform scanner missed the attack;
(b) SCS static regex detected dangerous capabilities;
(c) DMR disclosure matching confirmed the capabilities were undisclosed;
(d) TDM temporal analysis would flag the forward-drift injection.
This confirms the multi-layer design is necessary and
no single layer is sufficient.

\section{Detector Feature Ablation}
\label{app:ablation}

Table~\ref{tab:ablation} shows AUROC when each feature group is used
in isolation (Logistic Regression, 5-fold CV).

\begin{table}[h]
\centering
\caption{Feature group ablation for temporal drift detector.}
\label{tab:ablation}
\small
\begin{tabular}{@{}llr@{}}
\toprule
\textbf{Group} & \textbf{Key Features} & \textbf{AUROC} \\
\midrule
Drift core & CES delta, text diff, new kw & 0.979 \\
Disclosure & SDS, changelog len, vague & 0.979 \\
Family drift & Per-family CES deltas & 0.980 \\
\textbf{All combined} & & \textbf{1.000} \\
\bottomrule
\end{tabular}
\end{table}

\section{Responsible Disclosure}
\label{app:disclosure}

We notified both \clawhub{} and \skillssh{} of the T6 cases
\textit{[redacted for anonymous review]}
days before submission.
\clawhub{} acknowledged and confirmed all 40 cases.
\skillssh{} did not respond within the disclosure window.
We have withheld specific skill identifiers in this submission;
they will be released as supplementary material upon acceptance.

\section{LLM Annotation Protocol}
\label{app:llm_protocol}

\emph{Note: LLM-based semantic adjudication is deferred to future work.
This appendix describes the planned methodology.}

We propose to evaluate IC classification robustness using five independent
open-source LLM annotators on \NskillsDMR{} \clawhub{} skills.

\paragraph{Planned Models.}
We plan to deploy five fully open-weight models locally via vLLM~\cite{kwon2023vllm}
on an HPC cluster, with no commercial API dependence:
\begin{itemize}[topsep=2pt,itemsep=1pt]
\item \textbf{M1}: Qwen3-8B (Alibaba Cloud) --- chain-of-thought reasoning, multilingual
\item \textbf{M2}: Llama-3.1-8B-Instruct (Meta AI) --- instruction-tuned, North American research
\item \textbf{M3}: Mistral-7B-v0.3-Instruct (Mistral AI) --- efficient European-origin model
\item \textbf{M4}: Gemma-3-4B-IT (Google DeepMind) --- compact instruction-tuned, safety-focused
\item \textbf{M5}: GPT-OSS-20B (OpenAI) --- open-weight, 20B parameter research model
\end{itemize}

This five-family ensemble spans models from Alibaba, Meta, Mistral, Google, and OpenAI,
maximizing the diversity of training data, RLHF alignment procedures, and architectural
design choices to minimize the risk of systematic single-model extraction artifacts.

\paragraph{Prompt template.}
Each annotator receives a structured prompt containing:
(a)~the skill's user-visible description (truncated to 400 chars),
(b)~the developer's changelog (truncated to 500 chars),
(c)~statically detected new risk keywords,
(d)~a unified diff of the SKILL.md between versions (capped at 150 lines),
(e)~the CES drift score, and
(f)~the SDS disclosure score.

The system prompt instructs the model to classify the update as:
\texttt{CONFIRMED} (undisclosed dangerous capability),
\texttt{FALSE\_POSITIVE} (contextually benign or adequately disclosed), or
\texttt{UNCLEAR} (insufficient evidence).
The model must respond with structured JSON including a label,
one-sentence reason, specific risk evidence, and a boolean
\texttt{disclosed} flag.

\paragraph{Consensus.}
The final label uses majority vote ($\geq3/5$ agreement).
Each annotator operates independently with no access to
the other annotators' outputs.
We use the resulting consensus
to estimate review yield on sampled pipeline outputs, not to
construct a universal ground-truth benchmark.

\paragraph{Reproducibility.}
All prompts, raw outputs, and the annotation script
(\texttt{08b\_claude\_native\_m1.py}, \texttt{08c\_gpt\_native\_m2.py})
are included in our artifact release.
